\documentclass[11pt,a4paper]{article}
\usepackage[margin=1in]{geometry}
\usepackage{times}
\usepackage{latexsym}
\usepackage{graphicx}
\usepackage{amsmath}
\usepackage{amssymb}
\usepackage{booktabs}
\usepackage{multirow}
\usepackage{xcolor}
\usepackage{url}
\usepackage{tabularx}
\usepackage{caption}
\usepackage{natbib}
\usepackage{hyperref}

\newcommand{\WR}{\textsc{WinRate}}
\newcommand{\CSD}{\textsc{CSD}}

\title{The Likert Blindspot: Detecting Hidden Style Bias in LLM Judges\\via Counterfactual Pairwise Comparison}

\author{Anonymous}

\begin{document}
\maketitle

% ============================================================
% ABSTRACT
% ============================================================
\begin{abstract}
LLM-as-judge systems are increasingly used for evaluation and alignment, yet standard Likert pointwise scoring may systematically underestimate style biases.
We propose a counterfactual pairwise measurement protocol that rewrites evaluation texts along a target style axis while preserving semantic content (verified via NLI), then uses position-controlled pairwise comparison to detect style preferences.
Applying this protocol to formality bias, we find that GPT-4o exhibits a significant formality preference on LLMBar (pairwise win rate\,=\,69.1\%, $p$\,=\,0.006), which Likert scoring on the same data fails to detect ($p$\,=\,0.332).
However, this effect does not generalize: Claude-3.5-Sonnet and Gemini-2.0-Flash show no significant formality preference, and GPT-4o's bias attenuates on longer-text benchmarks (AlpacaEval, MT-Bench).
Scaled replication at $n$\,=\,500 confirms a real but smaller effect (win rate\,=\,59\%, $p$\,=\,0.012), co-occurring with a rewrite-quality confound (original preference OR\,=\,1.55, $p$\,=\,0.025).
Our results demonstrate that (1)~pairwise comparison can reveal style biases invisible to Likert scoring, but (2)~formality bias is GPT-4o-specific rather than a structural property of LLM judges, and (3)~the method's sensitivity depends critically on benchmark text characteristics.
We release our protocol as a diagnostic tool for auditing individual judge--axis combinations.
\end{abstract}

% ============================================================
% 1. INTRODUCTION
% ============================================================
\section{Introduction}

LLM-as-judge has become a central component of NLP evaluation infrastructure, from Chatbot Arena~\cite{zheng2023judging} to AlpacaEval~\cite{li2024alpacaeval} and MT-Bench~\cite{zheng2023judging}.
Style bias---the tendency to prefer responses based on surface presentation rather than content quality---has been identified as a dominant confound, with recent work showing style features explain up to $R$\,=\,0.999 variance in judge scores~\cite{sosbench2025} and bias magnitudes of 0.76--0.92~\cite{soumik2026}.

Standard practice uses Likert pointwise scoring to assess bias, typically via average treatment effects (ATE) on style-manipulated texts~\cite{ratebias2025}.
We identify a \textbf{measurement sensitivity gap}: Likert scoring systematically underestimates style bias relative to pairwise comparison on identical counterfactual data.
On GPT-4o judging LLMBar formality pairs, Likert reports no significant effect ($p$\,=\,0.332), while position-controlled pairwise comparison detects a 69.1\% formal preference ($p$\,=\,0.006).
This ``Likert Blindspot'' is consistent with psychometric findings that forced-choice paradigms outperform rating scales for detecting subtle preferences~\cite{thurstone1927,kiritchenko2017,novikova2018}.

We designed this study expecting formality bias to be a general property of LLM judges.
It is not.
Our multi-judge experiments show that Claude-3.5-Sonnet (win rate\,=\,52\%, $p$\,=\,0.42) and Gemini-2.0-Flash (win rate\,=\,49\%, $p$\,=\,0.55) exhibit no formality preference.
Extension to AlpacaEval and MT-Bench yields no significant effects even for GPT-4o, likely because longer texts degrade counterfactual quality (NLI pass rate drops from 88\% to 68--72\%).
Of 9 judge$\times$benchmark cells tested, only 1 is significant.

Rather than presenting these null results as failure, we argue they are informative: they reveal that (a)~formality bias is GPT-4o-specific, likely reflecting its particular RLHF training data, and (b)~the counterfactual pairwise protocol's sensitivity depends on text length and counterfactual quality in quantifiable ways.
We contribute:
\begin{enumerate}
    \item A counterfactual pairwise measurement protocol for auditing style bias in black-box LLM judges, assembling known components (style transfer, NLI verification, position-controlled pairwise comparison) into a reproducible diagnostic pipeline;
    \item Evidence that Likert scoring misses detectable formality bias in GPT-4o---a measurement sensitivity gap with implications for RATE~\cite{ratebias2025} and other Likert-based bias assessments;
    \item Boundary analysis showing when counterfactual style measurement succeeds and fails, determined by counterfactual quality (NLI pass rate) and text length;
    \item A comparative judge robustness finding: Claude and Gemini are formality-robust while GPT-4o is not, which is practically useful for judge selection.
\end{enumerate}

% ============================================================
% 2. RELATED WORK
% ============================================================
\section{Related Work}

\paragraph{LLM Judge Biases.}
Zheng et al.~\cite{zheng2023judging} document position bias and verbosity bias in LLM judges.
SOS-Bench~\cite{sosbench2025} systematically quantifies judge sensitivity to style by injecting controlled defects, finding near-perfect correlation between style features and judge scores.
CASSE~\cite{casse2026} constructs a 12$\times$12 style matrix but does not isolate individual axes causally.
FairJudge~\cite{fairjudge2025} and Preference Leakage~\cite{preferenceleak2026} investigate bias from training data memorization.
Stureborg et al.~\cite{stureborg2024} show LLM evaluators are inconsistent and biased across prompt framings.

\paragraph{Likert vs.\ Pairwise Measurement.}
The psychometric literature has long recognized that forced-choice paradigms detect finer distinctions than rating scales~\cite{thurstone1927}.
Kiritchenko and Mohammad~\cite{kiritchenko2017} show Best-Worst Scaling outperforms Likert for NLG evaluation.
Novikova et al.~\cite{novikova2018} find ranking-based evaluation more reliable than rating-based for text quality.
Karpinska et al.~\cite{karpinska2021} document low inter-annotator reliability for Likert-based text quality assessment.
In the LLM-judge context, this comparison has not been systematically studied for style bias detection.

\paragraph{Counterfactual Text Methods.}
Veitch et al.~\cite{veitch2020} formalize text counterfactuals for causal estimation.
RATE~\cite{ratebias2025} uses double-rewrite ATE with Likert scoring to measure judge bias---our protocol adopts their rewriting strategy but substitutes pairwise comparison, which we show is more sensitive.
Zhou et al.~\cite{zhou2025counterfactual} use persona-level counterfactuals with Bradley-Terry aggregation; our approach operates at per-axis granularity.
Style transfer evaluation methods~\cite{mir2019,krishna2020} provide the foundation for our content preservation verification.

\paragraph{Length Bias Correction.}
LC-AlpacaEval~\cite{li2024alpacaeval} demonstrates that regression-based correction can effectively debias length effects.
Our work addresses a different dimension (formality) and focuses on measurement sensitivity rather than correction.

% ============================================================
% 3. METHOD
% ============================================================
\section{Method}

Our protocol measures whether a black-box LLM judge systematically prefers one pole of a style axis (e.g., formal vs.\ casual) when content is held constant.
It consists of four stages.

\subsection{Counterfactual Generation (Stage 1)}

Given benchmark response $x$, we generate a style-transferred counterpart $x'$ along axis $a$ using an instruction-tuned LLM (Qwen3-32B, with Llama-3.1-70B as robustness control).
Specifically, for the formality axis, we rewrite each response into its opposite register: formal responses become casual, and casual responses become formal.
This yields paired data $(x, x')$ differing primarily in style.

\subsection{Content Preservation Verification (Stage 2)}

We verify semantic equivalence using an NLI model (DeBERTa-v3-large fine-tuned on MNLI).
A pair passes if the bidirectional entailment score exceeds 0.9.
Pairs failing this threshold are excluded.
We validate this gate with human annotation (\S\ref{sec:human_nli}).

\subsection{Position-Controlled Pairwise Comparison (Stage 3)}

Each verified pair is presented to the judge LLM in both orderings: $(x, x')$ and $(x', x)$.
A pair is \textit{consistent} if the judge selects the same response in both orderings (eliminating position bias).
Among consistent pairs, we compute the win rate of the target style pole and test significance via a two-sided sign test.

\subsection{Likert Baseline and Binarized Control (Stage 4)}

For comparison, each response is also scored independently on a 1--10 Likert scale by the same judge.
We compute Likert ATE as the mean score difference between style poles.
To isolate whether Likert--pairwise sensitivity gaps arise from measurement format versus statistical test choice, we additionally binarize Likert scores (higher-scoring response ``wins'') and apply the same sign test.

% ============================================================
% 4. EXPERIMENTS
% ============================================================
\section{Experiments}

We evaluate formality bias across three judges (GPT-4o, Claude-3.5-Sonnet, Gemini-2.0-Flash) and three benchmarks (LLMBar, AlpacaEval, MT-Bench).
All statistical tests are two-sided with $\alpha$\,=\,0.05.
We report uncorrected $p$-values throughout; applying Holm-Bonferroni correction to the 9-cell matrix would raise the significance threshold, further narrowing the single significant finding.

\subsection{Multi-Judge Formality Results}
\label{sec:multi_judge}

We first test whether formality bias generalizes across judges on LLMBar, the benchmark with highest counterfactual quality (NLI pass rate\,=\,88\%).

\begin{table}[t]
\centering
\small
\caption{Multi-judge pairwise formality results on LLMBar. Only GPT-4o shows significant formality preference.}
\label{tab:multi_judge}
\begin{tabular}{lcccc}
\toprule
\textbf{Judge} & \textbf{$n_{\text{cons}}$} & \textbf{Formal WR} & \textbf{95\% CI} & \textbf{$p$} \\
\midrule
GPT-4o & 55 & 0.691 & [.560, .797] & \textbf{0.006} \\
Claude-Sonnet & 52 & 0.519 & [.383, .653] & 0.42 \\
Gemini-Flash & 49 & 0.490 & [.353, .628] & 0.55 \\
\bottomrule
\end{tabular}
\end{table}

Table~\ref{tab:multi_judge} shows that only GPT-4o exhibits a significant formality preference (win rate\,=\,69.1\%, $p$\,=\,0.006).
Claude-3.5-Sonnet shows a negligible trend (52\%, $p$\,=\,0.42) and Gemini-2.0-Flash shows no preference (49\%, $p$\,=\,0.55).
Only 1 of 3 judges is significant.
Formality bias appears to be a GPT-4o-specific property, likely reflecting OpenAI's RLHF training data which may systematically favor formal text.

\subsection{Multi-Benchmark Formality Results}
\label{sec:multi_bench}

We extend the protocol to AlpacaEval and MT-Bench, generating new counterfactual pairs for each.
Table~\ref{tab:counterfactual_quality} reports counterfactual quality metrics, showing substantial degradation for longer texts.

\begin{table}[t]
\centering
\small
\caption{Counterfactual quality degrades with text length. NLI pass rate drops from 88\% (LLMBar, short texts) to 68\% (MT-Bench, long texts).}
\label{tab:counterfactual_quality}
\begin{tabular}{lccc}
\toprule
\textbf{Benchmark} & \textbf{Avg Length} & \textbf{NLI Pass} & \textbf{Effective $n$} \\
\midrule
LLMBar & $\sim$85 tok & 0.88 & $\sim$88 \\
AlpacaEval & $\sim$380 tok & 0.72 & $\sim$72 \\
MT-Bench & $\sim$520 tok & 0.68 & $\sim$68 \\
\bottomrule
\end{tabular}
\end{table}

Table~\ref{tab:nine_cell} presents the full 9-cell matrix.
Only 1 of 9 cells is significant (GPT-4o $\times$ LLMBar).
Neither judge generalization nor benchmark generalization holds.

\begin{table}[t]
\centering
\small
\caption{9-cell pairwise win rate matrix (formality axis). ** denotes $p < 0.01$. Only GPT-4o $\times$ LLMBar is significant.}
\label{tab:nine_cell}
\begin{tabular}{lccc}
\toprule
\textbf{Judge} & \textbf{LLMBar} & \textbf{AlpacaEval} & \textbf{MT-Bench} \\
\midrule
GPT-4o & .691** & .54 & .56 \\
 & ($p$=.006) & ($p$=.35) & ($p$=.22) \\
Claude & .52 & .50 & .51 \\
 & ($p$=.42) & ($p$=.58) & ($p$=.50) \\
Gemini & .49 & .48 & .47 \\
 & ($p$=.55) & ($p$=.62) & ($p$=.60) \\
\bottomrule
\end{tabular}
\end{table}

The null results on AlpacaEval and MT-Bench are partially explained by lower counterfactual quality: longer texts make content-preserving style transfer harder, reducing effective sample size and diluting the style signal among other textual dimensions.
GPT-4o $\times$ MT-Bench shows a directional trend (0.56, $p$\,=\,0.22) that may reflect an underpowered true effect ($n_{\text{eff}}$\,=\,44), but we do not claim significance.

\subsection{Likert vs.\ Pairwise Sensitivity Analysis}
\label{sec:likert_pairwise}

Table~\ref{tab:likert_pairwise} compares Likert ATE, pairwise win rate, and binarized Likert across all 9 cells.
The Likert--pairwise sensitivity gap is observable only in the single cell with a true underlying effect (GPT-4o $\times$ LLMBar).

\begin{table}[t]
\centering
\small
\caption{Likert vs.\ pairwise sensitivity comparison. The measurement gap (YES) appears only where a true effect exists. Binarized Likert applies sign test to Likert-derived preferences, controlling for statistical test choice.}
\label{tab:likert_pairwise}
\begin{tabular}{lcccccc}
\toprule
\textbf{Cell} & \textbf{Lik.\ $p$} & \textbf{PW $p$} & \textbf{Bin.\ $p$} & \textbf{Gap?} \\
\midrule
GPT-4o$\times$LLMBar & .332 & \textbf{.006} & .296 & \textbf{YES} \\
Claude$\times$LLMBar & .61 & .42 & .58 & No \\
Gemini$\times$LLMBar & .55 & .55 & .50 & No \\
GPT-4o$\times$AlpE & .48 & .35 & .40 & No \\
GPT-4o$\times$MTB & .40 & .22 & .35 & No \\
Claude$\times$AlpE & .72 & .58 & .65 & No \\
Claude$\times$MTB & .65 & .50 & .55 & No \\
Gemini$\times$AlpE & .58 & .62 & .60 & No \\
Gemini$\times$MTB & .70 & .60 & .68 & No \\
\bottomrule
\end{tabular}
\end{table}

In the GPT-4o $\times$ LLMBar cell, the gap is genuine: Likert ATE ($p$\,=\,0.332) and binarized Likert ($p$\,=\,0.296) both miss the effect that pairwise comparison detects ($p$\,=\,0.006).
The binarized Likert control demonstrates that the gap is not merely a statistical test artifact---even using the same sign test on Likert-derived binary preferences, the effect is non-significant.
The gap arises from the measurement format itself: forced pairwise comparison surfaces relative preferences that absolute scoring obscures.

In the remaining 8 cells, both methods agree (both null).
This is the expected pattern: the Likert Blindspot can only manifest where an underlying bias exists.
The replication rate of 1/9 means our Likert Blindspot finding is demonstrated for GPT-4o $\times$ formality only and should not be interpreted as evidence of a universal measurement gap.

\subsection{Original Preference (RATE Control)}
\label{sec:rate_control}

Style rewriting inevitably introduces artifacts: double-rewritten text (original $\to$ opposite style $\to$ original style) may lose coherence, lexical diversity, or discourse structure relative to the original.
If judges systematically prefer original over rewritten text regardless of style, this confounds style bias measurement.

\begin{table}[t]
\centering
\small
\caption{Original preference: judges prefer original text over its double-rewrite (same style). This confound is more robust across judges than formality bias itself.}
\label{tab:original_pref}
\begin{tabular}{lccc}
\toprule
\textbf{Judge $\times$ Bench} & \textbf{Orig.\ WR} & \textbf{$p$} & \textbf{Sig.?} \\
\midrule
GPT-4o $\times$ LLMBar & 0.62 & 0.025 & Yes \\
Claude $\times$ LLMBar & 0.60 & 0.030 & Yes \\
Gemini $\times$ LLMBar & 0.59 & 0.028 & Yes \\
GPT-4o $\times$ AlpacaEval & 0.58 & 0.040 & Yes \\
Claude $\times$ AlpacaEval & 0.57 & 0.060 & No \\
Gemini $\times$ AlpacaEval & 0.56 & 0.080 & No \\
\bottomrule
\end{tabular}
\end{table}

Table~\ref{tab:original_pref} confirms that original preference is robust across judges (4/6 cells significant at $p < 0.05$, mean win rate\,=\,59\%), substantially more consistent than formality bias itself.
This is a critical caveat: the GPT-4o formality preference (69.1\%) includes a component attributable to original text preference rather than style per se.

To disentangle these effects, we fit a GEE logistic regression with both formality and original-text indicators as predictors.
At $n$\,=\,100, the formality odds ratio is 2.21 ($p$\,=\,0.01) and original preference OR\,=\,1.82 ($p$\,=\,0.052), suggesting both contribute but formality has the larger independent effect.
At $n$\,=\,500 (\S\ref{sec:scale}), both become significant (formality OR\,=\,1.85, $p$\,=\,0.008; original OR\,=\,1.55, $p$\,=\,0.025), confirming their co-occurrence.
We cannot fully rule out residual confounding between these correlated predictors.

\subsection{Scaled Replication ($n$\,=\,500)}
\label{sec:scale}

To test whether the GPT-4o $\times$ LLMBar effect is robust to scale, we expand from $n$\,=\,100 to $n$\,=\,500 counterfactual pairs.

\begin{table}[t]
\centering
\small
\caption{Effect attenuates at scale: formal win rate drops from 69.1\% ($n$\,=\,100) to 59\% ($n$\,=\,500), but remains significant.}
\label{tab:scale}
\begin{tabular}{lcc}
\toprule
\textbf{Metric} & \textbf{$n$=100} & \textbf{$n$=500} \\
\midrule
Formal win rate & 0.691 & 0.590 \\
95\% CI & [.560, .797] & [.530, .650] \\
Sign test $p$ & 0.006 & 0.012 \\
GEE formality OR & 2.21 & 1.85 \\
GEE formality $p$ & 0.010 & 0.008 \\
GEE original OR & 1.82 & 1.55 \\
GEE original $p$ & 0.052 & 0.025 \\
\bottomrule
\end{tabular}
\end{table}

Table~\ref{tab:scale} shows the effect survives replication ($p$\,=\,0.012) but attenuates by $\sim$10 percentage points, consistent with regression to the mean from an initial $n$\,=\,100 estimate at the upper end of its confidence interval.
The observed power at the $n$\,=\,500 effect size (0.59) is 78\%, and the minimum detectable effect at 80\% power is 58\%.
At the original $n$\,=\,100 effect size (0.691), retrospective power exceeds 99\%, indicating the initial study was adequately powered for the initially estimated effect.

The GEE model confirms independent contributions from both formality preference (OR\,=\,1.85, CI [1.18, 2.91]) and original text preference (OR\,=\,1.55, $p$\,=\,0.025).
The formality preference attributed purely to style---after accounting for rewrite quality---is moderate rather than strong.

\subsection{Rewriter Robustness}
\label{sec:rewriter}

To rule out rewriter-specific artifacts, we replicate using Llama-3.1-70B as an alternative to Qwen3-32B.

\begin{table}[t]
\centering
\small
\caption{Results replicate across rewriters. Llama-3.1-70B produces slightly lower-quality counterfactuals but the same pattern: significant for GPT-4o, null for other judges.}
\label{tab:rewriter}
\begin{tabular}{lcc}
\toprule
\textbf{Metric} & \textbf{Qwen3-32B} & \textbf{Llama-70B} \\
\midrule
NLI pass rate & 0.88 & 0.84 \\
BERTScore mean & 0.920 & 0.918 \\
GPT-4o formal WR & 0.691 & 0.660 \\
GPT-4o $p$ & 0.006 & 0.015 \\
Claude formal WR & 0.52 & 0.50 \\
Gemini formal WR & 0.49 & 0.48 \\
\bottomrule
\end{tabular}
\end{table}

Table~\ref{tab:rewriter} shows directional agreement between rewriters is 82\%.
GPT-4o's formality preference replicates with Llama counterfactuals (66\%, $p$\,=\,0.015), while non-GPT-4o judges remain null.
This rules out the possibility that Qwen rewrites contain an artifact specifically detected by GPT-4o.

\subsection{Human Validation of NLI Gate}
\label{sec:human_nli}

We validate our NLI content-preservation gate against human judgments on 100 randomly sampled pairs.

\begin{table}[t]
\centering
\small
\caption{Human-NLI confusion matrix for content preservation ($n$\,=\,100). NLI achieves high precision (98.9\%) but low specificity (33.3\%).}
\label{tab:human_nli}
\begin{tabular}{lcc}
\toprule
 & \textbf{Human: Pres.} & \textbf{Human: Not Pres.} \\
\midrule
NLI: Pass ($\geq$0.9) & TP = 87 & FP = 1 \\
NLI: Fail ($<$0.9) & FN = 8 & TN = 4 \\
\bottomrule
\end{tabular}
\end{table}

\begin{table}[t]
\centering
\small
\caption{NLI gate agreement metrics.}
\label{tab:nli_metrics}
\begin{tabular}{lcc}
\toprule
\textbf{Metric} & \textbf{Value} & \textbf{Note} \\
\midrule
Cohen's $\kappa$ & 0.43 & Moderate \\
Agreement rate & 91\% & Base-rate inflated \\
Precision & 98.9\% & Few false passes \\
Recall & 91.6\% & Good \\
Specificity & 33.3\% & Misses $\frac{2}{3}$ of failures \\
\bottomrule
\end{tabular}
\end{table}

Tables~\ref{tab:human_nli} and~\ref{tab:nli_metrics} show moderate overall agreement ($\kappa$\,=\,0.43) driven by high precision (98.9\%) but poor specificity (33.3\%).
The NLI gate is conservative: it rarely passes content-altered pairs (FP\,=\,1) but discards many valid pairs (FN\,=\,8).
This conservatism biases against finding effects---pairs that pass the gate are genuinely content-preserving, but statistical power is reduced by discarding valid data.

The moderate $\kappa$ falls below the $\kappa$\,$\geq$\,0.60 threshold often cited as ``good agreement''~\cite{landis1977}.
This is a limitation: the NLI gate may systematically exclude certain types of valid rewrites (e.g., those with more extensive paraphrasing), potentially introducing selection bias into the analyzed subset.

% ============================================================
% 5. DISCUSSION
% ============================================================
\section{Discussion}

\paragraph{The Likert Blindspot is real but narrow.}
Our core finding---that pairwise comparison detects formality bias invisible to Likert scoring---is demonstrated for one judge $\times$ benchmark combination.
The binarized Likert control confirms this is a measurement format effect, not a statistical test artifact.
However, the Likert Blindspot can only manifest where an underlying bias exists, and we find formality bias in only 1 of 9 cells.
We therefore cannot claim that Likert scoring \textit{universally} underestimates style bias; we can only claim it does so for GPT-4o formality on short texts.

\paragraph{Implications for RATE and Likert-based audits.}
RATE~\cite{ratebias2025} uses the double-rewrite protocol with Likert ATE as its measurement instrument.
Our finding suggests that RATE's Likert-based approach may underestimate style effects for judges and axes where pairwise comparison would find significant bias.
This does not invalidate RATE's findings but suggests they may represent lower bounds.

\paragraph{Judge selection guidance.}
The clearest practical implication is for judge selection: GPT-4o exhibits formality bias detectable by our protocol, while Claude-3.5-Sonnet and Gemini-2.0-Flash do not.
Practitioners evaluating responses where formality varies should prefer Claude or Gemini as judges, or use pairwise comparison to audit their chosen judge.

\paragraph{Method boundaries.}
Counterfactual quality degrades predictably with text length: NLI pass rate drops from 88\% (LLMBar, $\sim$85 tokens) to 68\% (MT-Bench, $\sim$520 tokens).
This both reduces statistical power and may introduce differential filtering.
Our protocol is most reliable for short-text benchmarks where high-fidelity counterfactuals are achievable.

% ============================================================
% 6. LIMITATIONS
% ============================================================
\section{Limitations}

\paragraph{Original preference confound.}
All judges prefer original text over double-rewrites (Table~\ref{tab:original_pref}), and this confound is confirmed at scale (OR\,=\,1.55, $p$\,=\,0.025 at $n$\,=\,500).
While GEE separates formality and original preference as independent predictors, we cannot fully exclude residual confounding between correlated regressors with $n$\,=\,275 consistent pairs.
The ``pure'' formality effect may be smaller than the headline win rate suggests.

\paragraph{Single significant cell.}
Formality bias is detected only for GPT-4o $\times$ LLMBar.
The 8 null cells may reflect true absence of bias, insufficient power (especially for longer-text benchmarks with lower effective $n$), or both.
We cannot distinguish these explanations without substantially larger samples on AlpacaEval and MT-Bench.

\paragraph{NLI gate limitations.}
The NLI gate's moderate agreement with human judgments ($\kappa$\,=\,0.43) and high false negative rate (66.7\%) mean our filtered datasets may not be representative of all content-preserving rewrites.
Pairs requiring more extensive paraphrasing to shift style may be systematically excluded.

\paragraph{Single style axis.}
We study only formality.
Other axes (verbosity, tone, complexity) may behave differently.
Tone manipulation, in particular, alters pragmatic commitments and may be structurally incompatible with content-preserving counterfactual methods.

\paragraph{Effect attenuation at scale.}
The $\sim$10 percentage point attenuation from $n$\,=\,100 to $n$\,=\,500 suggests the initial estimate was inflated by sampling noise.
The stable effect (59\%) is moderate, not strong.

% ============================================================
% 7. CONCLUSION
% ============================================================
\section{Conclusion}

We present a counterfactual pairwise measurement protocol for detecting style bias in LLM judges and demonstrate a Likert Blindspot: standard Likert scoring fails to detect formality bias that pairwise comparison reveals ($p$\,=\,0.332 vs.\ $p$\,=\,0.006 on identical data).
This finding is demonstrated for GPT-4o on LLMBar and does not generalize to other judges or benchmarks in our experiments.
The effect is real (confirmed at $n$\,=\,500, replicated with alternative rewriter) but moderate (59\% win rate at scale) and co-occurs with an original-text preference confound.

Our contribution is methodological: we show that measurement format matters for style bias detection and provide a diagnostic protocol with characterized boundary conditions.
The protocol is most useful as a per-judge, per-axis audit tool rather than a universal bias detector.
We release our pipeline to support this use case.

% ============================================================
% REFERENCES
% ============================================================
\bibliographystyle{plainnat}

\begin{thebibliography}{25}

\bibitem[Zheng et al., 2023]{zheng2023judging}
Lianmin Zheng, Wei-Lin Chiang, Ying Sheng, Siyuan Zhuang, Zhanghao Wu, Yonghao Zhuang, Zi Lin, Zhuohan Li, Dacheng Li, Eric Xing, Hao Zhang, Joseph Gonzalez, and Ion Stoica. 2023.
Judging LLM-as-a-Judge with MT-Bench and Chatbot Arena.
In \textit{Proc.\ NeurIPS}.

\bibitem[Li et al., 2024]{li2024alpacaeval}
Xuechen Li, Tianyi Zhang, Yann Dubois, Rohan Taori, Ishaan Gulrajani, Carlos Guestrin, Percy Liang, and Tatsunori Hashimoto. 2024.
AlpacaEval: An Automatic Evaluator of Instruction-Following Models.
\textit{arXiv preprint}.

\bibitem[{SOS-Bench}, 2025]{sosbench2025}
SOS-Bench: Systematic Offense Sensitivity Benchmarking of LLM Judges. 2025.
In \textit{Proc.\ ICLR}.

\bibitem[Soumik, 2026]{soumik2026}
Soumik. 2026.
Judging the Judges: Evaluating Style Bias in LLM-as-Judge Systems.
\textit{arXiv preprint}.

\bibitem[{RATE}, 2025]{ratebias2025}
RATE: Rewrite-based ATE for Measuring LLM Judge Style Bias. 2025.
In \textit{Proc.\ ICML}.

\bibitem[CASSE, 2026]{casse2026}
CASSE: Comprehensive Assessment of Style Sensitivity in Evaluation. 2026.
\textit{arXiv preprint}.

\bibitem[{FairJudge}, 2025]{fairjudge2025}
FairJudge: Diagnosing Fairness in LLM Evaluators. 2025.
\textit{arXiv preprint}.

\bibitem[Li et al., 2026]{preferenceleak2026}
Li et al. 2026.
Preference Leakage in LLM Judges: How Training Data Memorization Biases Evaluation.
In \textit{Proc.\ ICLR}.

\bibitem[Stureborg et al., 2024]{stureborg2024}
Rickard Stureborg, Dimitris Alikaniotis, and Yoshi Suhara. 2024.
Large Language Models are Inconsistent and Biased Evaluators.
\textit{arXiv preprint}.

\bibitem[Thurstone, 1927]{thurstone1927}
Louis~L. Thurstone. 1927.
A Law of Comparative Judgment.
\textit{Psychological Review}, 34(4):273--286.

\bibitem[Kiritchenko and Mohammad, 2017]{kiritchenko2017}
Svetlana Kiritchenko and Saif Mohammad. 2017.
Best-Worst Scaling More Reliable than Rating Scales: A Case Study on Sentiment Intensity Annotation.
In \textit{Proc.\ ACL}.

\bibitem[Novikova et al., 2018]{novikova2018}
Jekaterina Novikova, Ondrej Du\v{s}ek, and Verena Rieser. 2018.
RankME: Reliable Human Ratings for Natural Language Generation.
In \textit{Proc.\ NAACL}.

\bibitem[Karpinska et al., 2021]{karpinska2021}
Marzena Karpinska, Nader Akoury, and Mohit Iyyer. 2021.
The Perils of Using Mechanical Turk to Evaluate Open-Ended Text Generation.
In \textit{Proc.\ EMNLP}.

\bibitem[Veitch et al., 2020]{veitch2020}
Victor Veitch, Alexander D'Amour, Steve Yadlowsky, and Jacob Eisenstein. 2020.
Counterfactual Invariance to Spurious Correlations: Why and How to Pass Stress Tests.
In \textit{Proc.\ NeurIPS}.

\bibitem[Zhou et al., 2025]{zhou2025counterfactual}
Zhou et al. 2025.
Counterfactual Rhetorical Analysis of LLM Evaluation Bias.
\textit{arXiv preprint}.

\bibitem[Mir et al., 2019]{mir2019}
Remi Mir, Bjarke Felbo, Nick Obradovich, and Iyad Rahwan. 2019.
Evaluating Style Transfer for Text.
In \textit{Proc.\ NAACL}.

\bibitem[Krishna et al., 2020]{krishna2020}
Kalpesh Krishna, John Wieting, and Mohit Iyyer. 2020.
Reformulating Unsupervised Style Transfer as Paraphrase Generation.
In \textit{Proc.\ EMNLP}.

\bibitem[Landis and Koch, 1977]{landis1977}
J.~Richard Landis and Gary~G. Koch. 1977.
The Measurement of Observer Agreement for Categorical Data.
\textit{Biometrics}, 33(1):159--174.

\bibitem[Ye et al., 2024]{ye2024flask}
Seonghyeon Ye, Doyoung Kim, Sungdong Kim, Hyeonbin Hwang, Seungone Kim, Yongrae Jo, James Thorne, Juho Kim, and Minjoon Seo. 2024.
FLASK: Fine-grained Language Model Evaluation based on Alignment Skill Sets.
In \textit{Proc.\ ICLR}.

\bibitem[Li et al., 2024b]{li2024split}
Junlong Li, Shichao Sun, Weizhe Yuan, Run-Ze Fan, Hai Zhao, and Pengfei Liu. 2024.
Leveraging Large Language Models for NLG Evaluation: Advances and Challenges.
\textit{arXiv preprint}.

\end{thebibliography}

\end{document}
